\section{x86-64 汇编语法总览}

\subsection{本章要解决的问题}

前面几章已经建立了三条基础线：

\begin{lstlisting}
C++ source -> compiler -> assembly text -> machine code bytes
object/type/layout -> address formula -> memory operand
instruction bytes -> CPU state transition -> performance behavior
\end{lstlisting}

现在正式进入 x86-64 汇编。很多初学者第一次看到汇编会有两个相反误区：一种是觉得汇编只是“更丑的 C”；另一种是觉得汇编完全不可读。两者都不对。汇编不是 C，它更接近 ISA 暴露给软件的机器动作；但汇编也不是随机符号，只要掌握语法、寄存器、操作数、指令宽度和 ABI 上下文，就可以逐条推导。

本章目标不是把所有 x86 指令一次性背完，而是建立读汇编的基本方法。你完成本章后，应该能做到：

\begin{itemize}
  \item 看懂一条 Intel 语法汇编指令的基本结构。
  \item 区分立即数、寄存器操作数、内存操作数。
  \item 识别通用寄存器及其 64/32/16/8 位别名。
  \item 理解 \code{byte ptr}、\code{dword ptr}、\code{qword ptr} 等宽度标记。
  \item 从 \code{[base + index*scale + displacement]} 还原 C++ 数组和结构体访问。
  \item 用“读哪些状态、写哪些状态”的方式注释汇编。
  \item 比较 \code{-O0} 和 \code{-O3} 汇编为什么差异巨大。
  \item 知道 Intel 语法和 AT\&T 语法的关键差异，避免读资料时混淆。
\end{itemize}

\begin{keyidea}
读汇编的核心不是把每条指令机械翻译成一行 C++，而是恢复数据流、控制流和地址公式。C++ 源码只是汇编的一种可能来源；真正可靠的是逐条说明指令读写了哪些寄存器、内存和 flags。
\end{keyidea}

\subsection{先生成一份你能控制的汇编}

学习汇编不能只看网上片段。你要从自己写的 C++ 生成汇编，知道编译命令、优化级别、目标语法和反汇编工具。先创建 \filepath{labs/ch10_assembly_syntax/minimal.cpp}：

\begin{lstlisting}[language=C++]
extern "C" int add3(int x, int y, int z)
{
    return x + y + z;
}

extern "C" int load_i32(int const* p, int i)
{
    return p[i];
}
\end{lstlisting}

生成 Intel 语法汇编：

\begin{lstlisting}[language=bash]
clang++ -std=c++20 -O3 -S -masm=intel minimal.cpp -o minimal.s
\end{lstlisting}

生成目标文件并反汇编：

\begin{lstlisting}[language=bash]
clang++ -std=c++20 -O3 -c minimal.cpp -o minimal.o
objdump -d -Mintel minimal.o > minimal.objdump.txt
\end{lstlisting}

如果你的系统主要使用 LLVM 工具链，也可以用：

\begin{lstlisting}[language=bash]
llvm-objdump -d --x86-asm-syntax=intel minimal.o > minimal.objdump.txt
\end{lstlisting}

\code{minimal.s} 是汇编器输入，通常包含标签、伪指令、调试/展开信息等。\code{minimal.objdump.txt} 是从目标文件机器码反汇编出来的文本，会显示机器码字节和指令地址。学习时两者都要看：

\begin{itemize}
  \item \code{.s} 更接近编译器输出，适合看函数边界和汇编器伪指令。
  \item \code{objdump} 更接近最终二进制，适合看机器码字节、真实指令地址和链接前符号。
\end{itemize}

\begin{warningbox}
不同编译器、版本、优化级别、目标 CPU、命令行参数都会生成不同汇编。本书的汇编是“合理示例”，不是要求你的机器逐字一致。学习重点是解释自己的输出。
\end{warningbox}

\subsection{一条汇编指令的基本结构}

Intel 语法里，一条普通指令大致长这样：

\begin{lstlisting}
mnemonic destination, source
\end{lstlisting}

例如：

\begin{lstlisting}
mov eax, edi
add eax, esi
ret
\end{lstlisting}

这段代码可以来自：

\begin{lstlisting}[language=C++]
extern "C" int add2(int x, int y)
{
    return x + y;
}
\end{lstlisting}

在 System V AMD64 ABI 下，前两个 \code{int} 参数在 \register{edi} 和 \register{esi}，返回值在 \register{eax}。逐条解释：

\begin{center}
\begin{tabular}{p{0.21\linewidth}p{0.33\linewidth}p{0.34\linewidth}}
\toprule
指令 & 读什么 & 写什么 \\
\midrule
\code{mov eax, edi} & \register{edi} & \register{eax}，同时清零 \register{rax} 高 32 位 \\
\code{add eax, esi} & \register{eax}、\register{esi} & \register{eax} 和部分 flags \\
\code{ret} & 栈顶返回地址、\register{rsp} & \register{rip}、\register{rsp} \\
\bottomrule
\end{tabular}
\end{center}

这里有一个初学阶段必须接受的事实：\instruction{add eax, esi} 的目的操作数既是输入也是输出。它不是“产生一个新变量”，而是把结果写回 \register{eax}：

\begin{lstlisting}
eax = eax + esi
\end{lstlisting}

因此读汇编时，你应该维护一个“当前每个寄存器代表什么”的笔记。比如：

\begin{lstlisting}
before:
    edi = x
    esi = y

mov eax, edi
    eax = x

add eax, esi
    eax = x + y

ret
    return eax
\end{lstlisting}

\subsection{操作数：立即数、寄存器和内存}

x86-64 指令的操作数常见三类：立即数、寄存器、内存。

\subsubsection{立即数}

\term{immediate}{立即数} 是编码在指令里的常量：

\begin{lstlisting}
mov eax, 42
add rdi, 16
and eax, 255
\end{lstlisting}

立即数不是内存地址，也不需要额外 load。它是机器码字节的一部分。比如 \code{add rdi, 16} 中的 \code{16} 会被编码进指令。

\subsubsection{寄存器操作数}

寄存器操作数直接读写 CPU 的架构寄存器：

\begin{lstlisting}
mov rax, rdi
add eax, esi
xor ecx, ecx
\end{lstlisting}

寄存器访问通常比内存访问快得多，但寄存器数量有限。编译器优化很大一部分工作，就是尽量把热数据保存在寄存器中，减少 load/store。

\subsubsection{内存操作数}

方括号表示内存操作数：

\begin{lstlisting}
mov eax, dword ptr [rdi]
mov eax, dword ptr [rdi + rsi*4]
mov qword ptr [rsp + 8], rax
\end{lstlisting}

方括号里的表达式计算地址；\code{byte ptr}、\code{dword ptr}、\code{qword ptr} 表示访问宽度。注意：

\begin{lstlisting}
mov eax, edi        ; register -> register
mov eax, [rdi]      ; memory at address rdi -> register
\end{lstlisting}

少一个方括号，语义完全不同。第一条把 \register{edi} 的值复制到 \register{eax}；第二条把 \register{rdi} 当作地址，从该地址读取内存。

\subsubsection{一般不能内存到内存}

多数普通 x86 指令不能同时有两个内存操作数。例如下面这种通常不合法：

\begin{lstlisting}
mov dword ptr [rdi], dword ptr [rsi]    ; invalid for ordinary mov
\end{lstlisting}

需要用寄存器中转：

\begin{lstlisting}
mov eax, dword ptr [rsi]
mov dword ptr [rdi], eax
\end{lstlisting}

这对性能分析很重要。源码里看起来像一次赋值，机器层面可能是一次 load 加一次 store。

\subsection{通用寄存器和别名}

x86-64 有 16 个常用通用整数寄存器。每个寄存器可以用不同宽度的名字访问低位部分：

\begin{center}
\begin{tabular}{lllll}
\toprule
64 位 & 32 位 & 16 位 & 8 位低位 & 常见 ABI 含义 \\
\midrule
\register{rax} & \register{eax} & \register{ax} & \register{al} & 返回值、临时 \\
\register{rbx} & \register{ebx} & \register{bx} & \register{bl} & callee-saved \\
\register{rcx} & \register{ecx} & \register{cx} & \register{cl} & 第 4 参数、移位计数 \\
\register{rdx} & \register{edx} & \register{dx} & \register{dl} & 第 3 参数、乘除辅助 \\
\register{rsi} & \register{esi} & \register{si} & \register{sil} & 第 2 参数 \\
\register{rdi} & \register{edi} & \register{di} & \register{dil} & 第 1 参数 \\
\register{rbp} & \register{ebp} & \register{bp} & \register{bpl} & 帧指针或 callee-saved \\
\register{rsp} & \register{esp} & \register{sp} & \register{spl} & 栈指针 \\
\bottomrule
\end{tabular}
\end{center}

\register{r8} 到 \register{r15} 也有类似名字：

\begin{center}
\begin{tabular}{llll}
\toprule
64 位 & 32 位 & 16 位 & 8 位 \\
\midrule
\register{r8} & \register{r8d} & \register{r8w} & \register{r8b} \\
\register{r9} & \register{r9d} & \register{r9w} & \register{r9b} \\
\register{r10} & \register{r10d} & \register{r10w} & \register{r10b} \\
\register{r11} & \register{r11d} & \register{r11w} & \register{r11b} \\
\register{r12} & \register{r12d} & \register{r12w} & \register{r12b} \\
\register{r13} & \register{r13d} & \register{r13w} & \register{r13b} \\
\register{r14} & \register{r14d} & \register{r14w} & \register{r14b} \\
\register{r15} & \register{r15d} & \register{r15w} & \register{r15b} \\
\bottomrule
\end{tabular}
\end{center}

\subsubsection{32 位写入会清零高 32 位}

x86-64 一个非常重要的规则是：写 32 位寄存器会把对应 64 位寄存器高 32 位清零。

\begin{lstlisting}
mov eax, edi
\end{lstlisting}

这条指令写 \register{eax}，也会让 \register{rax} 的高 32 位变成 0。可以这样理解：

\begin{lstlisting}
rax = zero_extend_64(edi)
\end{lstlisting}

但写 16 位或 8 位寄存器不会自动清零高位：

\begin{lstlisting}
mov ax, di      ; only low 16 bits are written
mov al, dil     ; only low 8 bits are written
\end{lstlisting}

因此部分寄存器写入可能保留旧值的高位。现代编译器通常更喜欢用 32 位写入来打破不必要的旧值依赖。

\begin{warningbox}
不要把 \register{eax}、\register{ax}、\register{al} 当成完全独立寄存器。它们是 \register{rax} 的不同宽度视图。读汇编时必须知道一条指令到底写了多少位。
\end{warningbox}

\subsection{操作数宽度：机器到底读写几个字节}

Intel 语法中，访问内存时常用宽度标记：

\begin{center}
\begin{tabular}{lll}
\toprule
标记 & 字节数 & 常见 C++ 类型或用途 \\
\midrule
\code{byte ptr} & 1 & \code{char}、\code{std::uint8_t} \\
\code{word ptr} & 2 & \code{std::uint16_t}、x86 word \\
\code{dword ptr} & 4 & \code{int}、\code{float}、\code{std::uint32_t} \\
\code{qword ptr} & 8 & \code{long}、指针、\code{double}、\code{std::uint64_t} \\
\code{xmmword ptr} & 16 & 128-bit SIMD memory operand \\
\code{ymmword ptr} & 32 & 256-bit AVX memory operand \\
\bottomrule
\end{tabular}
\end{center}

例如：

\begin{lstlisting}
movzx eax, byte ptr [rdi]
mov   eax, dword ptr [rdi]
mov   rax, qword ptr [rdi]
movss xmm0, dword ptr [rdi]
movsd xmm0, qword ptr [rdi]
\end{lstlisting}

这些指令即使地址相同，访问宽度和解释规则也不同：

\begin{itemize}
  \item \instruction{movzx eax, byte ptr [rdi]} 读取 1 字节并零扩展。
  \item \instruction{mov eax, dword ptr [rdi]} 读取 4 字节整数到 \register{eax}。
  \item \instruction{mov rax, qword ptr [rdi]} 读取 8 字节整数或指针。
  \item \instruction{movss xmm0, dword ptr [rdi]} 读取单精度浮点标量。
  \item \instruction{movsd xmm0, qword ptr [rdi]} 读取双精度浮点标量。
\end{itemize}

\begin{keyidea}
内存地址只说明从哪里开始，操作数宽度说明读写多少字节，指令语义和寄存器类型说明如何解释这些字节。三者缺一不可。
\end{keyidea}

\subsection{内存寻址模式：地址公式如何写进指令}

x86-64 最常见的内存操作数形式是：

\begin{lstlisting}
[base + index * scale + displacement]
\end{lstlisting}

其中 \code{scale} 只能是 1、2、4、8。不是所有部分都必须出现：

\begin{center}
\begin{tabular}{p{0.35\linewidth}p{0.48\linewidth}}
\toprule
形式 & 含义 \\
\midrule
\code{[rdi]} & 地址是 \register{rdi} \\
\code{[rdi + 8]} & 地址是 \register{rdi} 加 8 字节 \\
\code{[rdi + rsi*4]} & 地址是 \register{rdi} 加 \register{rsi} 乘 4 \\
\code{[rdi + rsi*4 + 12]} & 基址、索引缩放、常量偏移同时出现 \\
\code{[rip + symbol]} & RIP-relative，常用于全局对象或常量池 \\
\bottomrule
\end{tabular}
\end{center}

\subsubsection{案例：数组 load}

C++：

\begin{lstlisting}[language=C++]
extern "C" int load_i32(int const* p, int i)
{
    return p[i];
}
\end{lstlisting}

可能汇编：

\begin{lstlisting}
movsxd rsi, esi
mov    eax, dword ptr [rdi + rsi*4]
ret
\end{lstlisting}

假设执行前：

\begin{center}
\begin{tabular}{lll}
\toprule
寄存器 & 值 & 含义 \\
\midrule
\register{rdi} & \code{0x1000} & \code{p[0]} 地址 \\
\register{esi} & \code{3} & 下标 \code{i} \\
\bottomrule
\end{tabular}
\end{center}

执行：

\begin{lstlisting}
movsxd rsi, esi
    rsi = sign_extend_64(3)
        = 3

mov eax, dword ptr [rdi + rsi*4]
    effective_address = 0x1000 + 3 * 4
                      = 0x100c
    load 4 bytes from 0x100c..0x100f into eax
\end{lstlisting}

\subsubsection{案例：结构体字段}

C++：

\begin{lstlisting}[language=C++]
struct S {
    int a;
    double b;
    int c;
};

extern "C" int get_c(S const* p)
{
    return p->c;
}
\end{lstlisting}

常见布局：

\begin{center}
\begin{tabular}{llll}
\toprule
字节范围 & 内容 & 大小 & 说明 \\
\midrule
0..3 & \code{a} & 4 & 字段 \\
4..7 & padding & 4 & 让 \code{b} 8 字节对齐 \\
8..15 & \code{b} & 8 & 字段 \\
16..19 & \code{c} & 4 & 字段 \\
20..23 & tail padding & 4 & 让 \code{sizeof(S)} 为 8 的倍数 \\
\bottomrule
\end{tabular}
\end{center}

可能汇编：

\begin{lstlisting}
mov eax, dword ptr [rdi + 16]
ret
\end{lstlisting}

这里 \code{+ 16} 不是魔法数字，而是 \code{offset(S, c)}：

\begin{lstlisting}
addr(p->c) = addr(*p) + offset(S, c)
           = rdi      + 16
\end{lstlisting}

\subsubsection{案例：结构体数组字段}

C++：

\begin{lstlisting}[language=C++]
struct Pair {
    int a;
    int b;
};

extern "C" int get_b(Pair const* p, int i)
{
    return p[i].b;
}
\end{lstlisting}

公式：

\begin{lstlisting}
sizeof(Pair)    = 8
offset(Pair, b) = 4
addr(p[i].b)    = base + i * 8 + 4
\end{lstlisting}

可能汇编：

\begin{lstlisting}
movsxd rsi, esi
mov    eax, dword ptr [rdi + rsi*8 + 4]
ret
\end{lstlisting}

如果 \register{rdi} 是 \code{0x2000}，\register{rsi} 是 \code{5}：

\begin{lstlisting}
effective_address = 0x2000 + 5 * 8 + 4
                  = 0x202c
load 4 bytes from 0x202c..0x202f
\end{lstlisting}

\subsection{\instruction{lea}：看起来像内存，实际不访问内存}

\instruction{lea} 是学习 x86 汇编必须尽早掌握的指令。它的名字是 load effective address，但它只计算地址表达式，不读取内存。

\begin{lstlisting}
lea rax, [rdi + rsi*4 + 8]
\end{lstlisting}

语义是：

\begin{lstlisting}
rax = rdi + rsi * 4 + 8
\end{lstlisting}

不是：

\begin{lstlisting}
rax = memory[rdi + rsi * 4 + 8]
\end{lstlisting}

对比：

\begin{lstlisting}
lea rax, [rdi + rsi*4 + 8]        ; compute address-like integer
mov eax, dword ptr [rdi + rsi*4 + 8] ; load 4 bytes from memory
\end{lstlisting}

编译器经常用 \instruction{lea} 做整数乘法常数：

\begin{lstlisting}
lea eax, [rdi + rdi*2]      ; eax = x * 3
lea eax, [rdi + rdi*4]      ; eax = x * 5
lea eax, [rdi*8 + 16]       ; eax = x * 8 + 16
\end{lstlisting}

这解释了为什么你在非指针代码里也会看到 \instruction{lea}：

\begin{lstlisting}[language=C++]
extern "C" int mul5_add1(int x)
{
    return x * 5 + 1;
}
\end{lstlisting}

可能生成：

\begin{lstlisting}
lea eax, [rdi + rdi*4 + 1]
ret
\end{lstlisting}

\begin{mentalmodel}
看到 \instruction{lea} 时先问：它的结果后面被当作地址使用，还是只是整数计算结果？\instruction{lea} 本身不产生 cache miss，因为它不访问内存。
\end{mentalmodel}

\subsection{flags：有些指令会留下条件信息}

x86 有 RFLAGS/EFLAGS。很多整数指令会更新 flags，例如 \instruction{add}、\instruction{sub}、\instruction{cmp}、\instruction{test}。后续条件跳转或条件移动会读取这些 flags。

本章先只建立直觉：

\begin{center}
\begin{tabular}{lll}
\toprule
flag & 名称 & 常见用途 \\
\midrule
ZF & zero flag & 结果是否为 0，相等比较 \\
SF & sign flag & 结果最高位是否为 1 \\
CF & carry flag & 无符号进位/借位 \\
OF & overflow flag & 有符号溢出 \\
\bottomrule
\end{tabular}
\end{center}

例如：

\begin{lstlisting}
cmp edi, esi
setg al
\end{lstlisting}

\instruction{cmp edi, esi} 本质上计算 \code{edi - esi} 来设置 flags，但不保存减法结果。\instruction{setg al} 根据 flags 把 \register{al} 设置为 0 或 1。第 12 章会系统讲控制流和条件码；这里只要记住：读汇编时，有些数据依赖不在通用寄存器里，而在 flags 里。

\subsection{标签、函数和汇编器伪指令}

看一段 \code{clang++ -S} 生成的汇编，除了真正的机器指令，还会看到伪指令：

\begin{lstlisting}
    .text
    .globl  add2
    .type   add2,@function
add2:
    mov     eax, edi
    add     eax, esi
    ret
    .size   add2, .-add2
\end{lstlisting}

常见项：

\begin{center}
\begin{tabular}{p{0.25\linewidth}p{0.58\linewidth}}
\toprule
写法 & 含义 \\
\midrule
\code{.text} & 后面进入代码段 \\
\code{.globl add2} & 符号 \code{add2} 对链接器可见 \\
\code{.type add2,@function} & 标记符号类型是函数 \\
\code{add2:} & 标签，代表当前位置地址 \\
\code{.size add2, .-add2} & 记录函数大小 \\
\bottomrule
\end{tabular}
\end{center}

这些伪指令不是 CPU 执行的指令。它们给汇编器、链接器、调试器和异常展开机制使用。后续讲 ELF、链接、ABI 和 unwinding 时会展开。

\subsection{Intel 语法和 AT\&T 语法}

本书默认使用 Intel 语法，因为它和 Intel/AMD 手册、Compiler Explorer 的 Intel 输出、很多性能资料更接近。你仍然必须认识 AT\&T 语法，因为 GNU 工具链和很多历史资料会用它。

同一条指令：

\begin{lstlisting}
mov eax, dword ptr [rdi + rsi*4]      ; Intel
movl (%rdi,%rsi,4), %eax              ; AT&T
\end{lstlisting}

关键差异：

\begin{center}
\begin{tabular}{p{0.23\linewidth}p{0.31\linewidth}p{0.31\linewidth}}
\toprule
差异 & Intel & AT\&T \\
\midrule
操作数顺序 & destination, source & source, destination \\
寄存器写法 & \code{rax} & \code{\%rax} \\
立即数写法 & \code{42} & \code{\$42} \\
内存写法 & \code{[rdi + rsi*4]} & \code{(\%rdi,\%rsi,4)} \\
宽度表示 & \code{dword ptr} & 指令后缀 \code{movl} \\
\bottomrule
\end{tabular}
\end{center}

再看几个对应关系：

\begin{lstlisting}
Intel: mov rax, rdi
AT&T : movq %rdi, %rax

Intel: add eax, esi
AT&T : addl %esi, %eax

Intel: mov qword ptr [rsp + 8], rax
AT&T : movq %rax, 8(%rsp)
\end{lstlisting}

\begin{warningbox}
最容易错的是操作数顺序。Intel 是 \code{dst, src}，AT\&T 是 \code{src, dst}。读资料时先确认语法，否则会把数据流方向看反。
\end{warningbox}

\subsection{\code{-O0} 和 \code{-O3} 为什么差异巨大}

同一个 C++ 函数，在不同优化级别下可以完全不像同一个程序。看：

\begin{lstlisting}[language=C++]
extern "C" int f(int x)
{
    int y = x + 1;
    int z = y * 3;
    return z - 2;
}
\end{lstlisting}

\code{-O0} 的目标是容易调试，不是快。它可能生成栈帧，把局部变量放在栈上：

\begin{lstlisting}
push rbp
mov  rbp, rsp
mov  dword ptr [rbp - 4], edi
mov  eax, dword ptr [rbp - 4]
add  eax, 1
mov  dword ptr [rbp - 8], eax
mov  eax, dword ptr [rbp - 8]
imul eax, eax, 3
mov  dword ptr [rbp - 12], eax
mov  eax, dword ptr [rbp - 12]
sub  eax, 2
pop  rbp
ret
\end{lstlisting}

\code{-O3} 会做代数化简和寄存器分配，可能变成：

\begin{lstlisting}
lea eax, [rdi + rdi*2 + 1]
ret
\end{lstlisting}

因为：

\begin{lstlisting}
y = x + 1
z = y * 3 = 3*x + 3
return z - 2 = 3*x + 1
\end{lstlisting}

这个案例说明：

\begin{itemize}
  \item \code{-O0} 适合观察源代码结构、栈帧和调试信息。
  \item \code{-O3} 适合观察真实性能路径。
  \item 性能优化不能根据 \code{-O0} 汇编下结论。
  \item \code{-O3} 汇编可能已经和源码结构差很多，需要恢复语义而不是逐句对应。
\end{itemize}

\subsection{读汇编的固定流程}

以后每次读一个函数，按下面流程做笔记：

\begin{enumerate}
  \item 写出函数签名和 ABI 参数位置。
  \item 标出每个入口寄存器的含义，例如 \register{rdi} 是第一个指针，\register{esi} 是第二个 \code{int}。
  \item 从上到下维护寄存器含义表。
  \item 对每条内存操作写出有效地址公式。
  \item 标出访问宽度和读写字节范围。
  \item 标出哪些指令更新 flags，后面哪里读取 flags。
  \item 区分真正 load/store 和纯地址计算 \instruction{lea}。
  \item 最后再尝试恢复 C-like 伪代码。
\end{enumerate}

例如：

\begin{lstlisting}
movsxd rsi, esi
mov    eax, dword ptr [rdi + rsi*4 + 8]
add    eax, 1
ret
\end{lstlisting}

笔记应该像这样：

\begin{lstlisting}
entry:
    rdi = base pointer p
    esi = int index i

movsxd rsi, esi
    rsi = sign_extend_64(i)

mov eax, dword ptr [rdi + rsi*4 + 8]
    address = p + i * 4 + 8
    load width = 4 bytes
    possible source = p[i + 2] if p is int*

add eax, 1
    eax = loaded_value + 1

ret
    return eax
\end{lstlisting}

这种笔记比“这大概是数组访问”更有价值。高性能工程里，证据链必须精确到地址公式和访问宽度。

\subsection{案例：从 C++ 到汇编逐条解释}

创建函数：

\begin{lstlisting}[language=C++]
struct Item {
    int id;
    float score;
    int count;
};

extern "C" int update(Item* items, int i, int delta)
{
    items[i].count += delta;
    return items[i].count;
}
\end{lstlisting}

常见布局：

\begin{center}
\begin{tabular}{llll}
\toprule
字段 & offset & 大小 & 说明 \\
\midrule
\code{id} & 0 & 4 & \code{int} \\
\code{score} & 4 & 4 & \code{float} \\
\code{count} & 8 & 4 & \code{int} \\
\bottomrule
\end{tabular}
\end{center}

\code{sizeof(Item)} 通常是 12。因为 x86 scale 没有 12，编译器可能使用 \instruction{lea} 生成 \code{i * 3}，再乘 4：

\begin{lstlisting}
movsxd rax, esi
lea    rax, [rax + rax*2]
add    dword ptr [rdi + rax*4 + 8], edx
mov    eax, dword ptr [rdi + rax*4 + 8]
ret
\end{lstlisting}

逐条解释：

\begin{lstlisting}
entry:
    rdi = items
    esi = i
    edx = delta

movsxd rax, esi
    rax = sign_extend_64(i)

lea rax, [rax + rax*2]
    rax = i * 3

add dword ptr [rdi + rax*4 + 8], edx
    address = items + (i * 3) * 4 + 8
            = items + i * 12 + 8
            = addr(items[i].count)
    load old 4-byte count
    add delta
    store new 4-byte count
    update flags

mov eax, dword ptr [rdi + rax*4 + 8]
    reload 4-byte count into eax as return value

ret
\end{lstlisting}

为什么第二条 \instruction{mov} 要 reload？因为 \instruction{add dword ptr [mem], edx} 是 read-modify-write 内存指令，它把结果写回内存，但不把结果写入 \register{eax}。返回值需要在 \register{eax}，所以还要再 load。编译器也可能生成不同形式，比如先 load 到寄存器、加、store、再 return 寄存器。你要解释自己的输出。

\begin{deepdive}
read-modify-write 指令表面是一条汇编，但微架构内部通常涉及 load、ALU、store，并会占用 load/store 相关资源。后面讲 uops、ports、store buffer 和原子指令时，这一点会变得非常重要。
\end{deepdive}

\subsection{实验 10.1：建立第一份汇编阅读报告}

\begin{labbox}
创建 \filepath{labs/ch10_assembly_syntax/basic_reading.cpp}：

\begin{lstlisting}[language=C++]
#include <cstddef>
#include <cstdint>

extern "C" int add2(int x, int y)
{
    return x + y;
}

extern "C" int mix(int x)
{
    return x * 5 + 17;
}

extern "C" std::uint8_t load_u8(std::uint8_t const* p, std::size_t i)
{
    return p[i];
}

extern "C" int load_i32(int const* p, int i)
{
    return p[i + 2];
}
\end{lstlisting}

生成汇编和反汇编：

\begin{lstlisting}[language=bash]
clang++ -std=c++20 -O0 -S -masm=intel basic_reading.cpp -o basic_reading_O0.s
clang++ -std=c++20 -O3 -S -masm=intel basic_reading.cpp -o basic_reading_O3.s
clang++ -std=c++20 -O3 -c basic_reading.cpp -o basic_reading.o
objdump -d -Mintel basic_reading.o > basic_reading.objdump.txt
\end{lstlisting}

交付物：

\begin{enumerate}
  \item 对每个函数列出 ABI 参数寄存器和返回值寄存器。
  \item 从 \code{basic_reading_O3.s} 中摘出每个函数的核心指令。
  \item 对每条核心指令写“读什么、写什么、是否访问内存、是否更新 flags”。
  \item 对 \code{load_u8} 解释为什么可能出现 \instruction{movzx}。
  \item 对 \code{mix} 解释编译器是否使用 \instruction{lea}。
  \item 对 \code{load_i32} 写出地址公式，并说明 \code{+ 8} 从哪里来。
  \item 从 \code{objdump} 中记录机器码字节，说明汇编文本和机器码字节的关系。
\end{enumerate}
\end{labbox}

\subsection{实验 10.2：Intel 和 AT\&T 语法互译}

\begin{labbox}
用同一个源码生成两种语法：

\begin{lstlisting}[language=bash]
clang++ -std=c++20 -O3 -S -masm=intel basic_reading.cpp -o intel.s
clang++ -std=c++20 -O3 -S basic_reading.cpp -o att.s
\end{lstlisting}

交付物：

\begin{enumerate}
  \item 从两个文件中各选 8 条对应指令。
  \item 写成三列表：Intel、AT\&T、语义。
  \item 标出操作数顺序差异。
  \item 标出内存寻址写法差异。
  \item 标出宽度写法差异，例如 \code{dword ptr} 和 \code{movl}。
\end{enumerate}

报告中必须包含下面这种格式：

\begin{center}
\begin{tabular}{p{0.29\linewidth}p{0.29\linewidth}p{0.28\linewidth}}
\toprule
Intel & AT\&T & 语义 \\
\midrule
\code{mov eax, edi} & \code{movl \%edi, \%eax} & \code{eax = edi} \\
\code{add eax, esi} & \code{addl \%esi, \%eax} & \code{eax += esi} \\
\bottomrule
\end{tabular}
\end{center}
\end{labbox}

\subsection{实验 10.3：\code{-O0} 和 \code{-O3} 对比}

\begin{labbox}
创建 \filepath{labs/ch10_assembly_syntax/o0_o3_compare.cpp}：

\begin{lstlisting}[language=C++]
extern "C" int formula(int x)
{
    int a = x + 1;
    int b = a * 3;
    int c = b - x;
    return c + 7;
}
\end{lstlisting}

生成：

\begin{lstlisting}[language=bash]
clang++ -std=c++20 -O0 -S -masm=intel o0_o3_compare.cpp -o formula_O0.s
clang++ -std=c++20 -O3 -S -masm=intel o0_o3_compare.cpp -o formula_O3.s
\end{lstlisting}

交付物：

\begin{enumerate}
  \item 注释 \code{-O0} 中的栈帧建立、局部变量栈槽和返回路径。
  \item 注释 \code{-O3} 中的代数化简结果。
  \item 用数学式证明 \code{formula(x)} 可以化简成什么。
  \item 解释为什么不能用 \code{-O0} 汇编评估性能。
  \item 说明 \code{-O0} 对调试有什么价值。
\end{enumerate}
\end{labbox}

\subsection{本章作业}

\begin{exercisebox}
\subsubsection*{作业 1：逐条注释 20 条汇编}

从本章实验生成的 \code{-O3} 汇编中选择至少 20 条真实指令。每条指令必须写：

\begin{itemize}
  \item 指令文本。
  \item 操作数类型：立即数、寄存器、内存。
  \item 读取的寄存器或内存。
  \item 写入的寄存器或内存。
  \item 访问宽度。
  \item 是否更新 flags。
  \item 你认为它对应的 C++ 语义。
\end{itemize}

\subsubsection*{作业 2：地址公式恢复}

对下面汇编片段恢复可能的地址公式和 C-like 伪代码：

\begin{lstlisting}
movsxd rax, esi
mov    ecx, dword ptr [rdi + rax*4 + 12]
add    ecx, edx
mov    dword ptr [rdi + rax*4 + 12], ecx
mov    eax, ecx
ret
\end{lstlisting}

要求至少给出两种可能来源：

\begin{enumerate}
  \item \code{int*} 数组访问版本。
  \item 结构体数组字段访问版本。
\end{enumerate}

每种版本都必须解释 \code{+ 12} 的来源。

\subsubsection*{作业 3：语法互译}

把下面 Intel 语法翻译成 AT\&T 语法，并写出语义：

\begin{lstlisting}
mov rax, rdi
add eax, esi
movzx eax, byte ptr [rdi + rsi]
mov qword ptr [rsp + 8], rax
lea rax, [rdi + rsi*8 + 16]
\end{lstlisting}

再把下面 AT\&T 语法翻译成 Intel 语法：

\begin{lstlisting}
movl %edi, %eax
addl %esi, %eax
movzbl (%rdi,%rsi), %eax
movq %rax, 8(%rsp)
leaq 16(%rdi,%rsi,8), %rax
\end{lstlisting}

\subsubsection*{作业 4：机器码和指令长度}

用 \code{objdump} 记录至少 10 条指令的机器码字节。回答：

\begin{enumerate}
  \item 哪些指令是 1 字节？
  \item 哪些指令长度超过 4 字节？
  \item 立即数和 displacement 如何让指令变长？
  \item 为什么 x86 前端需要处理变长指令？
\end{enumerate}

\subsubsection*{作业 5：小型汇编阅读报告}

选择一个你自己写的 10 到 20 行 C++ 函数，要求包含：

\begin{itemize}
  \item 至少一个数组访问。
  \item 至少一个结构体字段访问。
  \item 至少一个整数表达式。
  \item 至少一个条件判断或比较。
\end{itemize}

生成 \code{-O0} 和 \code{-O3} 汇编。提交不少于 2000 字报告，必须包含：

\begin{itemize}
  \item 源码。
  \item 编译命令。
  \item 函数签名和 ABI 参数寄存器。
  \item \code{-O0} 和 \code{-O3} 汇编对比。
  \item 至少 15 条指令逐条注释。
  \item 所有内存访问的地址公式。
  \item 至少 5 条机器码字节记录。
  \item 你遇到的一个误判，以及如何纠正。
\end{itemize}
\end{exercisebox}

\subsection{验收标准}

完成本章后，你应该能做到：

\begin{itemize}
  \item 熟练区分 Intel 和 AT\&T 语法。
  \item 识别立即数、寄存器和内存操作数。
  \item 解释通用寄存器别名和 32 位写入清零高位。
  \item 从内存操作数写出有效地址公式。
  \item 解释 \code{byte ptr}、\code{dword ptr}、\code{qword ptr}、\code{xmmword ptr} 的访问宽度。
  \item 解释 \instruction{lea} 为什么不访问内存。
  \item 对简单函数生成 \code{-O0} 和 \code{-O3} 汇编，并逐条注释核心指令。
  \item 从反汇编中记录机器码字节，知道汇编文本最终对应指令编码。
  \item 不再把 \code{-O0} 汇编当作性能路径。
\end{itemize}

如果这些能力稳定，第 12 章的整数指令、第 13 章的控制流、第 14 章的 ABI、第 24 章的自动向量化和第 48 章的 AVX2 microkernel 才能真正读得进去。
